% Figure: algebraic identity rendered as a two-pan balance. Used to
% reinforce the balance archetype introduced at the end of section 1.1.
\begin{figure}[htbp]
\centering
\begin{tikzpicture}[
  pan/.style={draw=engblack, thick, fill=engblue!10,
    minimum width=3.2cm, minimum height=0.6cm},
  arm/.style={draw=engblack, thick},
  fulcrum/.style={draw=engblack, thick, fill=engblack},
  expr/.style={font=\small, align=center},
  labelnode/.style={font=\footnotesize\itshape, align=center},
]

% Beam
\draw[arm] (-4, 1.0) -- (4, 1.0);
% Suspension chains
\draw[arm] (-2.4, 1.0) -- (-2.4, 0.4);
\draw[arm] (-3.6, 1.0) -- (-3.6, 0.4);
\draw[arm] (2.4, 1.0) -- (2.4, 0.4);
\draw[arm] (3.6, 1.0) -- (3.6, 0.4);
% Pans
\node[pan] at (-3.0, 0.1) {};
\node[pan] at (3.0, 0.1) {};
% Pan contents
\node[expr] at (-3.0, 0.1) {$a^{2} - b^{2}$};
\node[expr] at (3.0, 0.1) {$(a-b)(a+b)$};
% Fulcrum
\draw[arm] (0, 1.0) -- (0, 1.8);
\node[fulcrum, regular polygon, regular polygon sides=3,
  shape border rotate=180, minimum size=0.5cm] at (0, 2.05) {};
\draw[arm] (-0.6, 2.4) -- (0.6, 2.4);

% Annotations
\node[labelnode] at (-3.0, -0.6) {one form\\of the expression};
\node[labelnode] at (3.0, -0.6) {an equivalent form\\of the same value};
\node[labelnode] at (0, 2.8) {identity\\(every legal substitution)};

\end{tikzpicture}
\caption[Algebraic identity as a balance]{An algebraic identity is a
balance. The two sides carry the same numerical value for every legal
substitution of the free variables. The difference-of-squares
identity is drawn here as the canonical case; the archetype recurs
across every later chapter of Volume II as conserved quantities,
equations of state, and exact identities preserved by linear
transformation.}
\label{fig:vol02:ch01:identity-balance}
\end{figure}
